% =============================================================
\section{Empirical results}\label{sec:results-empirical}
% =============================================================

This section answers four questions that a stimulation designer or
methodologist would put to the simulator. Does it agree with NEURON on
a realistic morphology before we trust any extracellular result? At a
given electrode distance and pulse width, what amplitude is needed to
evoke a spike? Where in the asymmetric biphasic waveform space does
the cell respond, given the same charge per pulse? And does the
infrastructure scale to design-space sweeps that NEURON cannot run? We
close with a gradient receipt that motivates the next step --- direct
optimisation of the waveform itself. All runs reported below are part
of the curated artifact set described in the reproducibility appendix;
wall-clock figures are on a single v5e chip unless stated otherwise.

\subsection{Validation against NEURON}\label{sec:res-neuron}

Before any stimulation result from this pipeline is trustworthy, the
simulator has to agree with NEURON on the same cell. We validate in
three steps, each adding one new mechanism on top of the previous so
that residuals can be attributed.

\paragraph{Setup and confounders.}
The NEURON reference uses the same BBP cADpyr229 morphology that
Jaxley loads, with three confounders equalised before measurement.
(i) Axon myelination is enabled by default in the BBP NEURON template;
we disable it for the comparison since our Jaxley implementation does
not model myelin. (ii) The \texttt{jaxley\_mech.channels.l5pc} ports
apply a fixed $Q_{10}$ temperature correction; we run NEURON at
$34\,^\circ$C to match. (iii) NEURON's \texttt{geom\_nseg} produces $1131$
segments on this morphology, while Jaxley with \texttt{ncomp}$\,=\,2$
and \texttt{max\_branch\_len}$\,=\,100$ produces $700$ compartments;
we report the mismatch rather than forcing them to match.

\paragraph{Step 1: channel kinetics.} We reduce the morphology to a
$20\,\mu$m$\times 20\,\mu$m cylinder (\texttt{parity\_single\_soma}),
so that the comparison exercises only channel kinetics and ion
reversals. With the full BBP Pyr and PV channel sets installed, spike
counts agree to $\pm 1$ over $\sim\!50$ spikes per cell (Pyr: $18$
NEURON / $16$ Jaxley; PV: $50$/$51$), and AP peak heights agree to
$\pm 0.04\,$mV across both cells. The match is consistent with a
source-level diff of the relevant \texttt{jaxley\_mech} channels
against the BBP \texttt{.mod} files.

\paragraph{Step 2: full-morphology Pyr, intracellular.} We instantiate
the full cADpyr229 morphology in both simulators and inject a
$0.5\,$nA somatic step current with extracellular coupling held out
(\texttt{parity\_bbp\_pyr}). Both simulators emit $3$ spikes
(\cref{fig:parity-intracellular}); first-spike latency differs by
$-50\,\mu$s ($2$ integration steps at $\Delta t = 0.025\,$ms), and AP
peak by $+0.49\,$mV. Subthreshold
waveform RMSE is $0.275\,$mV with Pearson $r = 0.9995$;
suprathreshold full-trace RMSE is $6.032\,$mV with $r = 0.9253$, where
the latter is dominated by spike-edge timing offsets. The
$-0.38\,$mV resting-potential residual is explained by our
Euclidean-distance approximation to BBP's path-distance apical $I_h$
gradient (\cref{sec:cells-bbp}); a path-distance traversal would close
it further.

\paragraph{Step 3: full-morphology Pyr, extracellular.} The headline
claim of this paper is that ECS works on a realistic branched cell.
We test it by comparing voltage traces with NEURON when both
simulators see the same point-source field on the same Pyr
morphology. After aligning Jaxley's compartment centres and NEURON's
segment readouts to the same physical sites
(\texttt{parity\_ecs.bbp\_pyr}), per-site voltage RMSE is
$0.387,\,0.553,\,0.505,\,0.322\,$mV at soma, apical, basal and axon
respectively (median $0.446\,$mV); the site-local $\phi$ matches
between the two simulators (\cref{fig:parity-ecs}). The ECS residuals should be read against
the Step-2 baseline, since the underlying solver and discretisation
differences carry over. Combined with the analytical
$O(\Delta x^2)$ convergence of \cref{sec:results-verify} on a uniform
straight cable, this anchors the activating-function pipeline at both
geometric extremes: closed-form correctness on a straight cable and
external parity on a branched real cell. The practical consequence is
that ECS-on-Pyr in this pipeline reproduces NEURON behaviour to
sub-millivolt residuals on a realistic morphology --- enough fidelity
to support design-space sweeps that NEURON cannot run in the same time
budget (\cref{sec:res-throughput}).

\begin{figure}[t]
  \centering
  \includegraphics[width=\linewidth]{figures/bbp_intracellular_parity.png}
  \caption{Step 2 (full-morphology Pyr, intracellular) cross-check.
  NEURON (black) and Jaxley (dashed) soma voltage traces under $0.1\,$nA
  subthreshold and $0.5\,$nA suprathreshold step currents on the same
  cADpyr229 morphology. Subthreshold RMSE $0.275\,$mV
  ($r = 0.9995$); suprathreshold RMSE $6.032\,$mV
  ($r = 0.9253$), dominated by spike-edge timing offsets
  attributable to the $700$-vs-$1131$ compartmentalisation difference.}
  \label{fig:parity-intracellular}
\end{figure}

\begin{figure}[t]
  \centering
  \includegraphics[width=\linewidth]{figures/bbp_ecs_parity.png}
  \caption{Step 3 (full-morphology Pyr, extracellular) cross-check.
  NEURON (black) and Jaxley (dashed) voltage traces at the soma and
  apical readout sites after aligning the two simulators' compartment
  centres to within $0\,\mu$m on the same Pyr morphology, plus per-site
  RMSE bars and the site-local $\phi$ that drives the ECS coupling.
  Median per-site RMSE $0.446\,$mV; the residuals carry over from the
  Step-2 baseline rather than reflecting a new ECS-specific
  discrepancy.}
  \label{fig:parity-ecs}
\end{figure}

\subsection{Strength--duration on the HH cable}\label{sec:res-sd}

At a given electrode distance and pulse width, what amplitude does it
take to evoke a spike? The strength--duration curve is the standard
first answer to that question, and underlies almost every clinical and
engineering decision downstream of it~\cite{4122088}. We define the threshold $A^*$ for a fixed
configuration $(d, \text{polarity}, t_p)$ as the smallest electrode
amplitude that elicits a single soma spike in the simulation:
\[
  A^*(d,\,\text{polarity},\,t_p)
  \;=\; \inf\bigl\{A \ge 0\,:\,F(w_{A,t_p};\,d) \text{ spikes}\bigr\}.
\]
Each $A^*$ is found by a one-dimensional binary search in $A$, with the
spike-detection criterion threshold-crossing of the soma voltage (see
\texttt{response.detect\_spike}). The full sweep is one binary search per
$(d,\,\text{polarity},\,t_p)$ triple --- four electrode distances
$\{25,\,50,\,100,\,200\}\,\mu$m, three waveform polarities (monophasic
cathodic, monophasic anodic, charge-balanced biphasic), and nine pulse
widths from $0.05$ to $2\,$ms --- for $108$ thresholds in total.

The search uses a uniform $[0,\,5000]\,\mu$A bracket with $14$
bisection iterations, giving a resolution of $\sim\!0.3\,\mu$A; spike
detection stays stable across this range for the HH cable. The whole
sweep runs in float-32. Cathodic and biphasic thresholds are well-resolved at
every $(d,\,t_p)$ pair and follow the textbook strength--duration
shape, declining toward a rheobase as $t_p$ grows
(\cref{fig:strength-duration}). Monophasic anodic elicits spikes only
at short pulse widths and short distances; it saturates at the bracket
ceiling for longer pulses. This polarity asymmetry is the discrete
analogue of the activating-function sign argument: cathodic stimulation
sets up a depolarising forcing in the central segments where the cable
fires first, whereas anodic stimulation depolarises the cable ends,
which only triggers under brief, geometrically-favourable
conditions~\cite{rattay1986analysis,Mcneal1976AnalysisOA}. For a
designer choosing parameters, the takeaway is that
ms-scale cathodic pulses minimise charge per spike at every distance
in the sweep, charge-balanced biphasic costs roughly two-fold more,
and pure anodic stimulation is only viable in the short-pulse corner.

\begin{figure}[t]
  \centering
  \includegraphics[width=\linewidth]{figures/strength_duration.png}
  \caption{Strength--duration curves on the $50$-compartment HH cable.
  One panel per electrode distance; one curve per waveform polarity.
  Threshold amplitude $|A^*|$ falls toward a rheobase as pulse width
  grows, and grows with electrode distance as the $1/d$ field decay
  implies. The dotted line is the binary-search ceiling ($5000\,\mu$A);
  pulse-width points missing from a curve indicate no spike below that
  ceiling, which characterises the anodic regime at long widths.}
  \label{fig:strength-duration}
\end{figure}

\subsection{Charge-balanced biphasic grid}\label{sec:res-grid}

Asymmetric biphasic pulses are the clinical lever for trading
activation selectivity against charge efficiency: at the same total
delivered charge per pulse, where on the $(T_p, T_n)$ plane does the
cell actually respond? A symmetric charge-balanced biphasic pulse has
two equal-and-opposite phases of equal duration; a clinical waveform
almost always has unequal phase durations (longer-low cathodic +
short-high anodic, or vice versa)~\cite{wongsarnpigoon2010energy}. To characterise
this region of waveform space, we hold the electrode at $d=100\,\mu$m
(the cleanest distance from the strength--duration sweep) and grid four
parameters: anodic phase duration $T_p$, cathodic phase duration $T_n$,
anodic amplitude $A_p$ (with cathodic amplitude $A_n$ pinned by
$A_p T_p = A_n T_n$ to keep total injected charge zero), and pulse-train
frequency $f$. We record spike count, latency, and peak voltage per
configuration rather than searching for a threshold, so each grid point
is one forward simulation.

The grid has $T_p,\,T_n \in \{50,\,100,\,200,\,500\}\,\mu$s
(four values each), $A_p \in \{10,\,20,\,30,\,40,\,80\}\,\mu$A, and
pulse-train frequencies $f \in \{100,\,200,\,500,\,1000\}\,$Hz, for
$315$ valid charge-balanced forward simulations; the five invalid
combinations have $T_p + T_n$ equal to the $1000\,$Hz period. The full grid
runs in $16\,$min on the v5litepod-4 pod-slice via
\texttt{vmap}-over-configuration with hierarchical scan-checkpointing
(\cref{sec:res-throughput}). Spike crossings concentrate along the
asymmetric $T_p = 500\,\mu$s edge at low-to-mid amplitudes
($A_p \le 30\,\mu$A), and at $A_p = 80\,\mu$A for $T_p = 200\,\mu$s
(\cref{fig:biphasic-grid}); the symmetric diagonal $T_p = T_n$ is far
less effective at the same charge per pulse. The
implication for waveform design is direct: asymmetry is not a marginal
parameter --- on this cell, at this distance, swapping a symmetric
biphasic pulse for an asymmetric one with the same charge can be the
difference between sub-threshold and reliable spiking.

\begin{figure}[t]
  \centering
  \includegraphics[width=\linewidth]{figures/biphasic_grid.png}
  \caption{Charge-balanced biphasic grid on the BBP L2/3 Pyr cell at
  $d = 100\,\mu$m. Each panel is a $(T_p, T_n)$ heatmap of the maximum
  soma voltage over the four pulse-train frequencies, with one panel
  per anodic amplitude $A_p$. The dotted diagonal $T_p = T_n$ marks the
  symmetric-biphasic line; cells where $\mathit{vmax}$ crosses
  $0\,$mV correspond to a soma spike at some frequency in the sweep.
  White cells indicate configurations for which all simulated
  frequencies returned non-finite voltage traces and therefore no
  interpretable $\mathit{vmax}$; these occur mainly in the
  long-$T_p$ high-amplitude regime. Five period-filling configurations
  ($T_p = T_n = 500\,\mu$s at $1000\,$Hz) are absent from the
  forward-simulation set.}
  \label{fig:biphasic-grid}
\end{figure}

\subsection{Throughput scaling on TPU}\label{sec:res-throughput}

The point of the simulator extension is to support design-space
sweeps that were previously infeasible. We measure how it scales
across two regimes that probe different bottlenecks. A
\emph{cheap-cells} sweep isolates per-cell compile and kernel-dispatch
overhead at small problem size: \texttt{ncomp}$=2$ on the BBP Pyr
($\sim\!700$ compartments per cell after branchpoint expansion), $B$
up to $1024$ on the v5litepod-4 pod-slice. A \emph{full-fidelity}
sweep is the actual scientific workload: $50$-compartment branches on
the same Pyr cell ($\approx\!17\,500$ compartments per cell), $B$ up
to $4096$ on the same pod-slice. Each $B$ value is a single
\texttt{jax.jit(jax.vmap(\ldots))} dispatch; we report the cached
(JIT-warm) timing.

The cheap-cells curve scales linearly to $B=1024$ at $5.1\,$ms per cell
($\sim\!195$ cells/s); the full-fidelity curve scales linearly through
$B=1024$, peaks at $7.28$ cells per second at $B=2048$, and rolls over
at $B=4096$ as the four chips approach arithmetic-intensity saturation
(\cref{fig:throughput-scaling}). The full-fidelity regime depends on
the sparse-$G$ and \texttt{checkpoint\_lengths} choices in the Method
(\cref{sec:method}); without either, the $(T,\,B,\,N)$ voltage
history overflows the v5e's $16\,$GB chip memory at $B \ge 256$.

\begin{figure}[t]
  \centering
  \includegraphics[width=\linewidth]{figures/throughput_v5p4_scaling.png}
  \caption{Jaxley batch throughput on the v5litepod-4 pod-slice for two
  workload regimes: cheap cells ($700$ compartments per cell, blue) and
  full-fidelity BBP Pyr ($\sim\!17\,500$ compartments per cell, red).
  Throughput scales linearly through $B=1024$ in both regimes; the
  full-fidelity curve peaks at $7.28$ cells per second at $B=2048$ and
  saturates at $B=4096$. JIT-warm timings.}
  \label{fig:throughput-scaling}
\end{figure}

\paragraph{Comparison with NEURON.} The same $700$-compartment Pyr cell
on a $T_\text{sim} = 200\,$ms protocol runs at $0.286$ simulations per
second under serial single-core NEURON (the default BBP setup), against
$9.67$ simulations per second for Jaxley at $B = 1000$ on a single TPU
once the JIT cache is warm --- a $\sim\!34\times$ wall-clock speedup
(\cref{fig:throughput}).
NEURON has parallel-CPU variants such as
CoreNEURON~\cite{Kumbhar2019CoreNEURONA} and
experimental GPU support that we did not benchmark; the relevant
comparison for our purposes is between a CPU-serial reference and an
accelerator-batched formulation. AI accelerators are SIMT/SIMD
throughput devices, and a batched-\texttt{vmap} formulation is the
workload pattern they are designed for; a CPU-bound serial framework
cannot expose the same batch parallelism without a similar
reformulation. The practical consequence: the parameter sweeps in
\cref{sec:res-sd,sec:res-grid} finish in tens of minutes, putting
sweeps over the full waveform-and-geometry design space within reach
of a single TPU job.

\begin{figure}[t]
  \centering
  \includegraphics[width=\linewidth]{figures/bbp_throughput.png}
  \caption{Wall-clock comparison on the same $700$-compartment Pyr cell
  and $T_\text{sim} = 200\,$ms protocol: serial single-core NEURON
  (blue) versus Jaxley on a single TPU chip (red). At $B=1000$, NEURON
  takes $58.3\,$min (measured), Jaxley takes $1.7\,$min ($103.4\,$ms
  per simulation, JIT-warm), giving a $\sim\!34\times$ wall-clock
  speedup. The dashed blue line is the linear extrapolation of the
  per-NEURON-sim cost.}
  \label{fig:throughput}
\end{figure}

\subsection{Gradient-based design signal}\label{sec:res-grad}

For the differentiability we built into the pipeline to matter for
stimulation design, the gradients have to point in the right
direction --- and their direction has to mean something physical.
Establishing that gradients exist (\cref{sec:results-verify}) is a
prerequisite for, but not equivalent to, showing that they carry usable
design signal. The position gradient reported in
\cref{sec:results-verify} admits a direct physical reading. With
the electrode at $(0,\,100,\,0)\,\mu$m, the negative $y$-component of
$\partial v_\mathrm{soma,\,peak} / \partial \mathbf{x}_e$ means moving
the electrode farther from the cell along its dorsal axis reduces the
soma response, consistent with the $1/d$ decay in the field model; the
horizontal components measure the local field anisotropy induced by the
asymmetric placement of dendrites relative to the electrode. Free-form
optimisation of $w(t)$ against a spike-elicitation target under
charge-balance and amplitude constraints is a natural next step and is
left as future work (\cref{sec:limitations}). What this section
establishes is that the same pipeline that runs the empirical sweeps
of \cref{sec:res-sd,sec:res-grid} also returns design-relevant
gradients in a single backward pass --- placing direct waveform
optimisation, rather than gradient-free search, within reach of the
same compute budget.

\begin{figure}[t]
  \centering
  \includegraphics[width=0.78\linewidth]{figures/gradient_arrow.png}
  \caption{Position-gradient direction on the BBP L2/3 Pyr cell. Grey
  points are the $700$-compartment morphology projected onto the
  $y$-$z$ plane; the soma is marked with a black dot, the electrode
  ($\mathbf{x}_e = (0,\,100,\,0)\,\mu$m) with a black cross. The
  orange arrow shows the direction of
  $\partial v_\mathrm{soma,\,peak} / \partial \mathbf{x}_e$, i.e.\ the
  direction in which moving the electrode increases the soma peak.}
  \label{fig:gradient-arrow}
\end{figure}
